\section{Phase 1}

\begin{frame}[plain]
  \vfill
  \begin{center}
    {\color{brandaccent}\Large\textbf{PHASE 1}}\\[0.4em]
    {\Large Chọn cách lấy đoạn văn}\\[1em]
    {\small\color{brandmuted}
     Trong \NumCorpusChunks{} đoạn, tìm \NumPoneRerankTopN{} đoạn nhiều khả năng chứa đáp án nhất.\\
     Không gọi LLM. Không tốn tiền API.}
  \end{center}
  \vfill
\end{frame}

\begin{frame}[shrink=20]{Thiết kế: giải đấu phân tầng}{23 cấu hình thay vì 72 tổ hợp}
  \small
  \begin{tabularx}{\textwidth}{@{}p{1.3cm} p{4.4cm} L@{}}
    \toprule
    \textbf{Vòng} & \textbf{Thử gì} & \textbf{Thắng} \\
    \midrule
    1 & 4 mô hình dense $\times$ 4 mô hình sparse & \win{BGE-M3 sparse} \\
    2 & 3 retriever $\times$ 3 reranker & \win{$+$ bge-reranker-large} \\
    3 & 3 kích thước chunk (256 / 512 / 1024) & \muted{512/64 --- không ai thắng rõ} \\
    \bottomrule
  \end{tabularx}

  \vspace{0.6em}
  Quán quân mỗi vòng mới được đi tiếp. Đây là đánh đổi có ý thức: rẻ hơn nhiều
  lần so với lưới đầy đủ, nhưng \textbf{không phát hiện được tương tác} giữa
  các lựa chọn ở hai vòng khác nhau.

  \vspace{0.4em}
  \textbf{Luật phán quyết khóa trước:} chênh lệch phải (1)~vượt biên độ nhiễu
  \NumNoiseFloor{}, \emph{và} (2)~có CI95 của hiệu theo cặp không chứa 0.

  \takeaway{Hai điều kiện, không phải một. Có ý nghĩa thống kê chưa đủ để đổi cấu hình.}
\end{frame}

\begin{frame}[shrink=20]{Vòng 1 --- Sparse áp đảo dense}
  \small
  \begin{tabularx}{\textwidth}{@{}L N N@{}}
    \toprule
    \textbf{So sánh} (nDCG@5, \texttt{resolved}) & \textbf{$\Delta$} & \textbf{CI95} \\
    \midrule
    BGE-M3 sparse \emph{vs} e5-base-v2 dense & \win{\NumSparseVsDense} & \multicolumn{1}{r}{\NumSparseVsDenseCI} \\
    BGE-M3 \emph{vs} BM25-okapi-stemmed & \NumBgeMThreeVsBmTwoFive & \multicolumn{1}{r}{\NumBgeMThreeVsBmTwoFiveCI} \\
    e5-base-v2 \emph{vs} bge-small-en-v1.5 & \NumDenseVsDense & \multicolumn{1}{r}{\NumDenseVsDenseCI} \\
    \bottomrule
  \end{tabularx}

  \vspace{0.8em}
  \begin{itemize}
    \item \textbf{Chỉ dòng đầu vượt được cả hai điều kiện.} Khoảng cách
          \NumSparseVsDense{} lớn hơn biên độ nhiễu \NumNoiseFloor{} và CI95 không chứa 0.
    \item Hai dòng sau: CI95 \textbf{chứa 0} $\rightarrow$ không phân định được.
          Không có ``mô hình dense tốt nhất'', và \emph{đó cũng là một kết quả}.
    \item \textbf{Vì sao sparse thắng:} EDA đã dự báo --- câu hỏi bám vào tên
          riêng, ngày tháng, con số. Khớp từ vựng bắt được đúng thứ đó; embedding
          làm nhòe nó đi.
  \end{itemize}

  \takeaway{Giả thuyết từ EDA, kiểm bằng bootstrap theo cặp, kết luận khớp dự báo.}
\end{frame}

\begin{frame}[shrink=20]{Vòng 2 và 3 --- Reranker cần, chunk size thì không}
  \small
  \textbf{Vòng 2 --- reranker thêm $+0{,}0659$ nDCG@5.}
  EDA đã đo trung vị \NumDistractorMedian{} đoạn đối thủ cùng chủ đề cho mỗi câu
  hỏi, nên việc lọc lại top~\NumPoneTopK{} xuống top~\NumPoneRerankTopN{} là bắt
  buộc. \muted{Hybrid dense+sparse: không chứng minh được có lợi.}

  \vspace{0.6em}
  \textbf{Vòng 3 --- kết quả null, và nó trung thực.}
  \begin{center}
    \begin{tabular}{@{}l r@{}}
      \toprule
      Chunk 512 / 64 & \NumChunkFiveOneTwoNdcg \\
      Chunk 1024 / 128 & \NumChunkOneZeroTwoFourNdcg \\
      Chunk 256 / 32 & \NumChunkTwoFiveSixNdcg \\
      \midrule
      Khoảng cách cao nhất -- thấp nhất & \NumChunkSpread \\
      \bottomrule
    \end{tabular}
  \end{center}

  \NumChunkSpread{} nằm \textbf{dưới} biên độ nhiễu \NumNoiseFloor{}, và cả ba
  khoảng CI95 chồng lấn nhau. Chọn 512/64 vì \textbf{lý do vận hành}, không
  phải vì nó thắng.

  \takeaway{Báo cáo một kết quả null đúng như nó là, thay vì tô nó thành chiến thắng.}
\end{frame}

\begin{frame}[shrink=20]{Ablation --- Contextual chunking}{Ba vòng đầu đổi \emph{retriever}. Vòng này đổi \emph{cách biểu diễn dữ liệu}.}
  \small
  Mọi chỉ mục ở Vòng 1--3 đều dựng từ \texttt{chunk["text"]} và không gì khác. Nên
  ``sparse thắng dense'' thực ra là ``sparse thắng dense \emph{trên cách biểu diễn
  này}''.

  \vspace{0.4em}
  \textbf{Giả thuyết:} chunk tiếp nối đến tay retriever mà không mang dấu hiệu nào
  cho biết nó thuộc bài nào.
  \begin{center}\footnotesize\ttfamily
  \begin{tabular}{@{}l l@{}}
    chunk 0 & "LOS ANGELES (CNN) -- Natalie Cole's search for a kidney ended..." \\
    chunk 1 & "...she said. The transplant was performed Tuesday..." \quad\rmfamily thuộc bài nào? \\
  \end{tabular}
  \end{center}

  \textbf{Can thiệp:} nối \NumCtxChars{} ký tự đầu bài báo vào trước mỗi chunk tiếp nối.
  \begin{center}\footnotesize
    \NumCorpusChunks{} chunk $-$ \NumCorpusArticles{} chunk đầu bài
    $=$ \textbf{\NumCtxContextualised{}} $=$ đúng số chunk được sửa (51,4\%)
  \end{center}

  \vspace{0.3em}
  \begin{tabularx}{\textwidth}{@{}L N N@{}}
    \toprule
    nDCG@5 (\texttt{resolved}, \NumCtxNSamples{} câu) & \textbf{plain} & \textbf{contextual} \\
    \midrule
    Dense (e5-base-v2) & \NumCtxDensePlain & \win{\NumCtxDenseCtx} \\
    Sparse (BGE-M3) & \win{\NumCtxSparsePlain} & \NumCtxSparseCtx \\
    \bottomrule
  \end{tabularx}

  \takeaway{Dense \NumCtxEffectDense{}, sparse \NumCtxEffectSparse{} --- hai hiệu ứng ngược chiều.}
\end{frame}

\begin{frame}[shrink=20]{Ablation --- đọc kết quả cho đúng}
  \small
  \begin{itemize}
    \item Dense \pass{\NumCtxEffectDense} \muted{\NumCtxEffectDenseCI} ---
          ngữ cảnh cho embedding thứ nó thiếu: biết chunk này thuộc bài nào.
    \item Sparse \fail{\NumCtxEffectSparse} \muted{\NumCtxEffectSparseCI} ---
          chính \NumCtxChars{} ký tự đó lặp trên mọi chunk cùng bài, \textbf{làm
          loãng tần suất từ hiếm}.
    \item Khoảng cách sparse--dense thu hẹp \NumCtxGapPlain{} $\rightarrow$
          \NumCtxGapCtx{}, nhưng \textbf{không đảo ngược}.
  \end{itemize}

  \vspace{0.3em}
  \textbf{Và kết luận này mạnh hơn bảng Vòng 1, chứ không yếu hơn:} nhánh dense ở
  đây dùng \textbf{tìm kiếm chính xác} thay vì HNSW --- vừa tái lập được (bỏ đúng
  nguồn nhiễu làm số dense trôi $0{,}0143$), vừa \textbf{công bằng hơn với dense}
  vì không mất recall do xấp xỉ. Sparse vẫn thắng trong điều kiện đó.

  \vspace{0.3em}
  \textbf{Hai chỗ phải nói rõ:}
  \begin{itemize}
    \item Đây là \emph{contextual chunking}, \textbf{không phải} \emph{metadata
          indexing}. Dataset không có metadata thật --- \texttt{url}, \texttt{author},
          \texttt{publish\_date} rỗng; \texttt{title} là \NumCtxChars{} ký tự đầu
          của chính body (kiểm 200/200 bài).
    \item CI ở đây bootstrap theo \emph{câu hỏi}, chưa gom cụm theo bài báo ---
          nhất quán với Phase 1, nhưng lỏng hơn chuẩn Phase 2.
  \end{itemize}

  \takeaway{Cả hai hiệu ứng đều dưới ngưỡng \NumNoiseFloor{} --- đo được, nhưng không đổi quyết định.}
\end{frame}

\begin{frame}[shrink=20]{Cấu hình khóa}
  \small
  \begin{center}
    BGE-M3 learned sparse, \texttt{top\_k=\NumPoneTopK}\quad$\rightarrow$\quad
    \texttt{BAAI/bge-reranker-large}, top~\NumPoneRerankTopN\quad$\rightarrow$\quad
    chunk 512/64
  \end{center}

  \vspace{0.6em}
  \begin{center}
    \begin{tabular}{@{}c c c c c@{}}
      \toprule
      \textbf{Hit@1} & \textbf{Hit@3} & \textbf{Hit@5} & \textbf{nDCG@5} & \textbf{Latency P50} \\
      \midrule
      \NumPoneHitOne & \NumPoneHitThree & \win{\NumPoneHitFive} & \NumPoneNdcgFive & \NumPoneLatency{} ms \\
      \bottomrule
    \end{tabular}
  \end{center}

  \vspace{0.6em}
  Chỉ mục được đóng băng và băm (\texttt{\NumLockedIndexSha\ldots}). Toàn bộ
  Phase 2 chạy lại đúng vết truy xuất này --- \textbf{không truy xuất lại lần nào}.

  \takeaway{Hit@5 $=$ \NumPoneHitFive{} sẽ trở thành \emph{trần} của Phase 2. Giữ con số này trong đầu.}
\end{frame}
